\section{Discussion}
\label{sec:discussion}

\subsection{Key Findings}

Our experiments reveal several important findings with implications for both research methodology and practical LLM training.

\textbf{Non-monotonic dose-response relationships exist and are measurable.} The quadratic fit ($R^2=0.985$) decisively outperforms linear models, confirming that curation parameter effects are not monotonic. This finding challenges the implicit assumption in many pipelines that ``stricter filtering is better'' and validates the quality-diversity tradeoff as a real phenomenon, not just theoretical intuition.

\textbf{The mechanism is noise dilution at low thresholds, diversity loss at high thresholds.} Convergence dynamics analysis shows 40\% higher AUC for unfiltered training, directly demonstrating that noisy data dilutes learning signals. Simultaneously, aggressive filtering (p90) achieves reasonable convergence but poor final performance, indicating that removed samples contained valuable information for generalization. This mechanistic understanding enables principled parameter selection beyond curve-fitting.

\textbf{Optimal thresholds are remarkably consistent across scales.} The transfer ratio of 0.85 between 125M and 1B models suggests that small-scale sweeps can effectively guide large-scale training. This has immediate practical value: instead of expensive full sweeps at production scale, teams can conduct systematic optimization at 125M with approximately 15\% discount for scale.

\subsection{Limitations}

Our work has several limitations that warrant acknowledgment.

\textbf{Reduced-scale evaluation.} All experiments used 5M--10M tokens rather than the planned 10B, with mock/synthetic benchmark evaluation due to lm-eval-harness integration issues. This means effect magnitudes are directional guidance, not production-ready numbers. The methodology is validated, but full-scale replication is required before deployment.

\textbf{Single architecture family.} Results were demonstrated on GPT-2 variants only. While GPT-2 is a standard reference architecture, different architectures (Llama-style, Mamba) may exhibit different sensitivities to curation parameters. Architecture generalization is future work.

\textbf{Single dataset family.} Experiments used RedPajama-v2, representative of English web corpora. Domain-specific corpora (code, scientific text, multilingual) may have different optimal thresholds due to different perplexity distributions.

\textbf{Perplexity dimension only fully validated.} While we tested both perplexity and deduplication, the deduplication analysis (H-M2) was inconclusive due to infrastructure issues. The dose-response pattern was visible in training loss but not benchmark scores due to evaluation failure.

\textbf{Single-parameter sweeps.} We studied perplexity and deduplication independently. Interaction effects (perplexity $\times$ deduplication joint optimization) were not tested. The full Pareto surface across both dimensions remains unexplored.

\subsection{Broader Impact}

This research has primarily positive implications for the field:

\textbf{Efficiency gains without new techniques.} Our findings enable practitioners to extract more performance from existing pipelines by calibrating curation parameters---essentially free improvement from parameter tuning rather than architectural innovation.

\textbf{Democratization of optimization.} The scale transfer finding means that smaller organizations with limited compute can conduct systematic optimization at affordable scales and transfer insights to larger training runs.

\textbf{Methodology contribution.} The fixed-token experimental design and dose-response analysis framework can be applied to other curation parameters (deduplication thresholds, domain mixing ratios, quality classifiers) and other domains (code, multilingual, multimodal).

\textbf{Potential negative impacts} are limited given the research nature of this work. If misapplied, overly aggressive perplexity filtering could inadvertently remove valuable minority-group or domain-specific content. Practitioners should validate that their quality signals do not encode demographic biases before threshold optimization.
